% Part: first-order-logic
% Chapter: completeness
% Section: identity

\documentclass[../../../include/open-logic-section]{subfiles}

\begin{document}

\olfileid{fol}{com}{ide}
\olsection{Identiteit}

\begin{explain}
De constructie van het termmodel uit de vorige paragraaf volstaat om de
volledigheid van de predicatenlogica vast te stellen voor
verzamelingen~$\Gamma$ die geen~$\eq$ bevatten. Het termmodel vervult
elke $!A \in \Gamma^*$ die geen~$\eq$ bevat (en dus elke
$!A \in \Gamma$). Het werkt echter niet wanneer~$\eq$ voorkomt. Dan kan
$\Gamma^*$ namelijk !!a{sentence}~$\eq[t][t']$ bevatten, terwijl in het
termmodel de waarde van elke term die term zelf is. Als $t$ en $t'$
verschillende termen zijn, zijn hun waarden in het termmodel---namelijk
respectievelijk $t$ en $t'$---dus verschillend, zodat $\eq[t][t']$
onwaar is. We kunnen dit verhelpen met een constructie die
``quotiëntvorming'' heet.
\end{explain}

\begin{defn}
  Laat $\Gamma^*$ een consistente en !!{complete} verzameling zinnen in
  $\Lang L$ zijn. We definiëren de relatie $\approx$ op de verzameling
  gesloten termen van~$\Lang L$ door
  \[
  t \approx t' \text{\quad dan en slechts dan als \quad} \eq[t][t'] \in \Gamma^*
  \]
\end{defn}

\begin{prop}
\ollabel{prop:approx-equiv}
De relatie $\approx$ heeft de volgende eigenschappen:
\begin{enumerate}
\item $\approx$ is reflexief.
\item $\approx$ is symmetrisch.
\item  $\approx$ is transitief.
\item Als $t \approx t'$, $f$ !!a{function} is en $t_1$, \dots,
  $t_{i-1}$, $t_{i+1}$, \dots, $t_n$ gesloten termen zijn, dan
\[
\Atom{f}{t_1,\dots, t_{i-1}, t, t_{i+1}, \dots, t_n} \approx
\Atom{f}{t_1,\dots, t_{i-1}, t', t_{i+1}, \dots, t_n}.
\]
\item Als $t \approx t'$, $R$ !!a{predicate} is en $t_1$, \dots,
  $t_{i-1}$, $t_{i+1}$, \dots, $t_n$ gesloten termen zijn, dan
\begin{multline*}
\Atom{R}{t_1,\dots, t_{i-1}, t, t_{i+1}, \dots, t_n} \in \Gamma^* \text{ dan en slechts dan als } \\
\Atom{R}{t_1,\dots, t_{i-1}, t', t_{i+1}, \dots, t_n} \in \Gamma^*.
\end{multline*}
\end{enumerate}
\end{prop}

\begin{proof}
Omdat $\Gamma^*$ consistent en !!{complete} is, geldt
$\eq[t][t'] \in \Gamma^*$ dan en slechts dan als
$\Gamma^* \Proves \eq[t][t']$. Het volstaat dus het volgende aan te tonen:
\begin{enumerate}
\item $\Gamma^* \Proves \eq[t][t]$ voor alle gesloten termen~$t$.
\item Als $\Gamma^* \Proves \eq[t][t']$, dan $\Gamma^* \Proves \eq[t'][t]$.
\item Als $\Gamma^* \Proves \eq[t][t']$ en $\Gamma^* \Proves
  \eq[t'][t'']$, dan $\Gamma^* \Proves \eq[t][t'']$.
\item Als $\Gamma^* \Proves \eq[t][t']$, dan
\[
\Gamma^* \Proves
\eq[\Atom{f}{t_1,\dots,t_{i-1},t,t_{i+1},,\dots,t_n}][\Atom{f}{t_1,\dots,t_{i-1},t',t_{i+1},\dots,t_n}]
\]
voor elke $n$-plaatsige !!{function}~$f$ en gesloten termen $t_1$, \dots,
$t_{i-1}$, $t_{i+1}$, \dots,~$t_n$.
\item Als $\Gamma^* \Proves \eq[t][t']$ en
$\Gamma^* \Proves
\Atom{R}{t_1,\dots,t_{i-1},t,t_{i+1},\dots,t_n}$, then
$\Gamma^* \Proves \Atom{R}{t_1,\dots,t_{i-1},t',t_{i+1},\dots,t_n}$
voor elk $n$-plaatsig !!{predicate}~$R$ en gesloten termen $t_1$, \dots,
$t_{i-1}$, $t_{i+1}$, \dots,~$t_n$.
\end{enumerate}
\end{proof}

\begin{prob}
Voltooi het bewijs van~\olref[fol][com][ide]{prop:approx-equiv}.
\end{prob}

\begin{defn}
Veronderstel dat $\Gamma^*$ een consistente en !!{complete} verzameling
in een taal~$\Lang L$ is, dat $t$ een gesloten term is en dat $\approx$
de relatie uit de vorige definitie is. Dan:
\[
\equivrep{t}{\approx} = \Setabs{t'}{t'\in \Trm[L], t \approx t'}
\]
and $\equivclass{\Trm[L]}{\approx} = \Setabs{\equivrep{t}{\approx}}{t \in \Trm[L]}$.
\end{defn}

\begin{defn}
\ollabel{defn:term-model-factor}
Laat $\Struct M = \Struct M(\Gamma^*)$ het termmodel voor~$\Gamma^*$ uit
\olref[mod]{defn:termmodel} zijn. Dan is
$\Struct{\equivclass{M}{\approx}}$ de volgende !!{structure}:
\begin{enumerate}
\item $\Domain{\equivclass{M}{\approx}} = \equivclass{\Trm[L]}{\approx}$.
\item $\Assign{c}{\equivclass{M}{\approx}} = \equivrep{c}{\approx}$
\item $\Assign{f}{\equivclass{M}{\approx}}(\equivrep{t_1}{\approx}, \dots,
  \equivrep{t_n}{\approx}) = \equivrep{\Atom{f}{t_1,\dots, t_n}}{\approx}$
\item $\tuple{\equivrep{t_1}{\approx}, \dots, \equivrep{t_n}{\approx}} \in
  \Assign{R}{\equivclass{M}{\approx}}$ dan en slechts dan als
  $\Sat{M}{\Atom{R}{t_1,\dots, t_n}}$, dat wil zeggen dan en slechts dan
  als $\Atom{R}{t_1,\dots,
  t_n} \in \Gamma^*$.
\end{enumerate}
\end{defn}

\begin{explain}
Merk op dat we $\Assign{f}{\equivclass{M}{\approx}}$ en
$\Assign{R}{\equivclass{M}{\approx}}$ voor elementen van
$\equivclass{\Trm[L]}{\approx}$ hebben gedefinieerd door ze te schrijven
als $\equivrep{t}{\approx}$, dus via \emph{representanten}~$t \in
\equivrep{t}{\approx}$. We moeten nagaan dat deze definities niet afhangen
van de keuze van die representanten. Andere keuzes~$t'$ die dezelfde
equivalentieklassen bepalen
($\equivrep{t}{\approx} = \equivrep{t'}{\approx}$), moeten dus hetzelfde
resultaat opleveren. Als $R$ bijvoorbeeld een eenplaatsig !!{predicate}
is, zegt de laatste clausule van de definitie dat
$\equivrep{t}{\approx} \in \Assign{R}{\equivclass{M}{\approx}}$ dan en
slechts dan als $\Sat{M}{\Atom{R}{t}}$. Als voor een andere term~$t'$ met
$t \approx t'$ daarentegen $\Sat/{M}{\Atom{R}{t}}$, dan zou de definitie
vereisen dat $\equivrep{t'}{\approx} \notin
\Assign{R}{\equivclass{M}{\approx}}$. Uit $t \approx t'$ volgt echter
$\equivrep{t}{\approx} = \equivrep{t'}{\approx}$, terwijl niet zowel
$\equivrep{t}{\approx} \in \Assign{R}{\equivclass{M}{\approx}}$ als
$\equivrep{t}{\approx} \notin \Assign{R}{\equivclass{M}{\approx}}$ kan
gelden. Volgens \olref{prop:approx-equiv} kan deze situatie dan ook niet
optreden.
\end{explain}

\begin{prop}
$\Struct{\equivclass{M}{\approx}}$ is goed gedefinieerd. Als $t_1$,
  \dots, $t_n$, $t_1'$, \dots, $t_n'$ gesloten termen zijn en
  $t_i \approx t_i'$, dan
\begin{enumerate}
\item $\equivrep{\Atom{f}{t_1,\dots, t_n}}{\approx} =
    \equivrep{\Atom{f}{t_1',\dots, t_n'}}{\approx}$, dat wil zeggen
  \[
  \Atom{f}{t_1,\dots, t_n} \approx \Atom{f}{t_1',\dots, t_n'}
  \]
  en
\item $\Sat{M}{\Atom{R}{t_1,\dots, t_n}}$ dan en slechts dan als
  $\Sat{M}{\Atom{R}{t_1',\dots, t_n'}}$, dat wil zeggen
  \[
    \Atom{R}{t_1,\dots, t_n} \in \Gamma^* \text{ dan en slechts dan als }
    \Atom{R}{t_1',\dots, t_n'} \in \Gamma^*.
  \]
\end{enumerate}
\end{prop}

\begin{proof}
Dit volgt uit \olref{prop:approx-equiv} door inductie naar~$n$.
\end{proof}

Net als bij het termmodel hebben we het volgende lemma nodig voordat we het
waarheidslemma bewijzen.

\begin{lem} 
\ollabel{lem:val-in-termmodel-factored} Laat $\Struct M = \Struct M
 (\Gamma^*)$. Dan geldt $\Value{t}{\equivclass{M}{\approx}} = \equivrep{t}
 {\approx}$.
\end{lem}
\begin{proof} 
Het bewijs is analoog aan dat van \olref[mod]{lem:val-in-termmodel}.
\end{proof}

\begin{prob}
Voltooi het bewijs van~\olref[fol][com][ide]{lem:val-in-termmodel-factored}.
\end{prob}

\begin{lem}
\ollabel{lem:truth}
$\Sat{\equivclass{M}{\approx}}{!A}$ dan en slechts dan als
$!A \in \Gamma^*$ voor alle zinnen~$!A$.
\end{lem}

\begin{proof}
Door inductie naar~$!A$, net als in het bewijs van
\olref[mod]{lem:truth}. Het enige geval dat extra aandacht vereist, is
$!A \ident \eq[t][t']$.
\begin{align*}
\Sat{\equivclass{M}{\approx}}{\eq[t][t']} & \text{ dan en slechts dan als } \equivrep{t}{\approx} = \equivrep{t'}{\approx}
\text{ (volgens de definitie van $\Struct{\equivclass{M}{\approx}}$)}\\
& \text{ dan en slechts dan als } t \approx t' \text{ (volgens de definitie van $\equivrep{t}{\approx}$)}\\
& \text{ dan en slechts dan als } \eq[t][t'] \in \Gamma^* \text{ (volgens de definitie van $\approx$).}
\end{align*}
\end{proof}

\begin{digress}
Hoewel $\Struct{M(\Gamma^*)}$ altijd !!{enumerable} en oneindig is, kan
$\Struct{\equivclass{M}{\approx}}$ eindig zijn: er kunnen slechts eindig
veel klassen $\equivrep{t}{\approx}$ zijn. Dat is te verwachten, want
$\Gamma$ kan !!{sentence}s bevatten die eisen dat elke !!{structure}
waarin ze waar zijn eindig is. Zo is
$\lforall[x][\lforall[y][x = y]]$ een consistente !!{sentence}, maar zij
wordt alleen vervuld in !!{structure}s met !!a{domain} dat precies één
!!{element} bevat.
\end{digress}

\end{document}
